% Einheit 2 von 5 – Bernoulli-Formel (bank/binomialverteilung/e2.jsonl, zusammenbau v0.1)
%% TODO Zeitmarke Einheit 2: Marken-Zeile nicht lesbar
%% TODO Zweigzeile Teil 1: was hier gelernt wird (Satz in Schülersprache aus dem Einheitstitel) – Einheit 2
\einheitenkopf[e2]{Einheit 2 von 5 · Bernoulli-Formel}
\zweigzeile{Abitur GK}
\setcounter{aufgabe}{12}

%% TODO Titel: Ich-kann-Satz fehlt in der Bank (Platzhalter: Kettenname) – Nr. 13
%% TODO Anweisung: ein Satz über der ersten Teilaufgabe fehlt in der Bank – Nr. 13
\begin{aufgabe}{Bernoulli-Formel}
%% TODO \rechenplatz{n}: Schrittzahl des Grundfalls fehlt in der Bank (nur bei fester Schreibform, 2.3 b)
\begin{teile}
\teil Ein Glücksrad wird 10-mal gedreht; X ist die Anzahl der Drehungen mit Rot. „Genau viermal Rot“ – welcher Ansatz passt? \kreuz{$P(X = 4)$} \\ \kreuz{$P(X \le 4)$} \\ \kreuz{$1 - P(X \le 3)$}
\teil Ein Glücksrad bleibt mit der Wahrscheinlichkeit $0{,}35$ auf Gold stehen und wird 6-mal gedreht; X ist die Anzahl der Drehungen mit Gold. Gib den Term für $P(X = 2)$ an und berechne ihn.
\teil 12\,\% der Smartphones einer Serie haben einen Displayfehler. 8 Geräte werden zufällig geprüft. Genau 3 mit Displayfehler – mit welcher Wahrscheinlichkeit? Runde auf vier Stellen. P ≈ \leerfeld
\teil Ein Versandhändler weiß, dass 8\,\% seiner Pakete zu spät ankommen. 30 Pakete werden zufällig ausgewählt. Mit welcher Wahrscheinlichkeit kommt keines davon zu spät an? Runde auf vier Stellen \hfill (Abitur 2017 GK). \leerfeld
\teil 90\,\% der Fluggäste erscheinen zu ihrem Flug. 20 Gäste haben gebucht. Mit welcher Wahrscheinlichkeit erscheinen genau 18? Zähle die Gäste, die nicht erscheinen, und runde auf vier Stellen. \leerfeld
\teil Ein Würfel wird viermal geworfen. Gib einen Term an, mit dem man die Wahrscheinlichkeit für „höchstens eine Sechs“ berechnen kann \hfill (Abitur 2026 GK).
\teil X ist binomialverteilt mit der Trefferwahrscheinlichkeit $\frac{2}{5}$. Vervollständige: $P(X = \square) = (\square \text{ über } 4) \cdot (\square)^3 \cdot \left(\frac{2}{5}\right)^4$ \hfill (Abitur 2021 GK).
\teil In einer Großstadt fahren 35\,\% der Pendler mit dem Rad. 120 Pendler werden zufällig befragt. Mit welcher Wahrscheinlichkeit fahren genau 40 davon mit dem Rad? Runde auf vier Stellen. \leerfeld
\teil Ein Glücksrad wird 20-mal gedreht; im Mittel bringen 7 Drehungen einen Gewinn. Mit welcher Wahrscheinlichkeit gibt es genau 7 Gewinne? Runde auf vier Stellen. \leerfeld
\teil 12\,\% der Autos auf einem Parkplatz sind Elektroautos. 50 Autos werden zufällig ausgewählt; X ist die Anzahl der Elektroautos. Die Tabelle zeigt Werte der summierten Binomialverteilung. Lies ab und berechne die Wahrscheinlichkeiten für genau 6 Elektroautos und für weniger als 6 Elektroautos \hfill (Abitur 2022 GK).

\sachtabelle{lcccc}{$k$ & 4 & 5 & 6 & 7}{$P(X \le k)$ & 0{,}2680 & 0{,}4353 & 0{,}6065 & 0{,}7533}
\end{teile}
\end{aufgabe}

%% TODO Titel: Ich-kann-Satz fehlt in der Bank (Platzhalter: Kettenname) – Nr. 14
%% TODO Anweisung: ein Satz über der ersten Teilaufgabe fehlt in der Bank – Nr. 14
\begin{aufgabe}{Bernoulli-Formel – Fehler finden, Begründen, Anwendung, Darstellungswechsel}
\begin{teile}
\teil Jonas berechnet die Wahrscheinlichkeit für genau zwei Sechsen bei fünf Würfen: \rechnung{P &= \left(\tfrac{1}{6}\right)^2 \cdot \left(\tfrac{5}{6}\right)^3 \approx 0{,}0161} Finde den Fehler.
\teil Bei drei Versuchen mit Treffer T und Niete N gibt es die Pfade TTN, TNT und NTT. Begründe, warum sie alle dieselbe Wahrscheinlichkeit haben.
\teil Eine Fluggesellschaft verkauft 52 Tickets für ein Flugzeug mit 50 Plätzen, weil erfahrungsgemäß 4\,\% der Gäste nicht erscheinen. Mit welcher Wahrscheinlichkeit reichen die Plätze nicht, erscheint also höchstens einer der Gäste nicht? Runde auf vier Stellen.
\teil Beschreibe ein Zufallsexperiment und ein Ereignis, dessen Wahrscheinlichkeit der Term $(6 \text{ über } 2) \cdot 0{,}15^2 \cdot 0{,}85^4$ angibt.
\end{teile}
\end{aufgabe}
